% !PRINTED 652
\weberpage{652}

Bei \gesperrt{Abel} findet sich für $k_0$ ein etwas anderer Ausdruck,
nämlich
\[
k_0 = C + B\sqrt{1+e^2} + \sqrt{h\left(1+e^2+\sqrt{1+e^2}\right)},
\]
der aus (12) hervorgeht, wenn man $a = 1:e$ setzt und dann $B$ und $h$
durch $Be$ und $he^2$ ersetzt. Der Ausdruck (12) ist also insofern
allgemeiner, als er auch den besonderen Fall $a=0$ umfasst, in dem die
biquadratische Gleichung in zwei quadratische Gleichungen zerfällt.

Um nun die Wurzel einer metacyklischen Gleichung 5\textsuperscript{ten}
Grades darzustellen, sind diese Ausdrücke für $k_0, k_1, k_2, k_3$ in
die Formel (9) des \S.\,186 zu substituiren. Nehmen wir $g=2$ an, so
werden die Exponenten $r_0, r_1, r_2, r_3$ der Reihe nach congruent mit
1, 2, 4, 8, also gleich 1, 2, 4, 3, und es ergiebt sich, wenn, wie in
\S.\,185
\[
\begin{aligned}
\tau_0 &= \sqrt[5]{k_0}, & \tau_1 &= \sqrt[5]{k_1} \\
\tau_2 &= \sqrt[5]{k_2}, & \tau_3 &= \sqrt[5]{k_3}
\end{aligned}
\]
gesetzt wird,
\weq{13}{
\begin{aligned}
\xi = A &+ K_0\, \tau_0^3\tau_1^4\tau_2^2\tau_3 + K_1\, \tau_1^3\tau_2^4\tau_3^2\tau_0 \\
        &+ K_2\, \tau_2^3\tau_3^4\tau_0^2\tau_1 + K_3\, \tau_3^3\tau_0^4\tau_1^2\tau_2.
\end{aligned}
}
Die Coëfficienten in diesem Ausdruck, $K_0, K_1, K_2, K_3$, sind
rationale Functionen der $k_0, k_1, k_2, k_3$, die durch cyklische
Permutation der $k$ selbst cyklisch permutirt werden. Nach \gesperrt{Abel}
ist
\[
K_0 = A_1 + A_2 k_0 + A_3 k_2 + A_4 k_0 k_2
\]
zu setzen, worin $A_1, A_2, A_3, A_4$ rational sind, was aber nicht
allgemein genug ist, weil zwischen $k_0 k_2$, $k_0$ und $k_2$ eine aus
(12) leicht abzuleitende lineare Relation besteht. Am einfachsten
drückt man die $k$ durch die drei Radicale
\[
\begin{aligned}
r &= \sqrt{1+a^2}, & \varrho &= \sqrt{h\left(1+a^2+a\sqrt{1+a^2}\right)} \\
\varrho' &= \sqrt{h\left(1+a^2-a\sqrt{1+a^2}\right)}, & \varrho\varrho' &= hr
\end{aligned}
\]
aus.

Man kann das Radical $\varrho'$ linear ausdrücken durch $r\varrho$ und
$\varrho$, wie man aus der Formel
\[
hr^2\varrho' = r\varrho\varrho'^2
\]
ersieht, wenn man rechts für $\varrho'^2$ seinen Ausdruck durch $r$
einsetzt und bedenkt, dass $r^2$ rational ist. So erhält man, wenn man
statt der cyklischen Permutation der $k$ die Vertauschung
\[
\begin{pmatrix}
r, & \varrho, & \varrho' \\
-r, & \varrho', & -\varrho
\end{pmatrix}
\]
%CONT>>
